\chapter{Как изучают структуру}\label{ch:methods}

Всё, что модель берёт на вход и с чем сравнивает результат, кем-то измерено: параметры решёток и несоответствие
гидрида --- рентгеном, текстура --- рентгеном или EBSD, линии растворимости --- калориметрией, ориентация и длина
гидридов --- на травленых шлифах. Чтобы понимать, насколько можно доверять этим числам, надо понимать, как они
получены. Глава идёт от рентгеновской трубки к шлифу.

\section{Рентгеновские лучи}

Рентгеновское излучение --- электромагнитные волны длиной $\sim0.01$--1~нм, сравнимой с межатомными расстояниями.
Именно поэтому кристалл для него --- дифракционная решётка.

В \textbf{рентгеновской трубке} электроны, ускоренные напряжением $V$ (десятки кВ), бьют в металлический анод. Получается
два вида излучения (рис.~\ref{fig:m1_spectrum}).
\begin{itemize}
\item \textbf{Тормозное} --- непрерывный спектр от торможения электронов. Его коротковолновая граница
  $\lambda_{\min} = hc/eV$ (весь электрон отдал энергию одному кванту): при 30~кВ это 0.041~нм.
\item \textbf{Характеристическое} --- узкие линии: электрон выбивает электрон с внутренней K-оболочки атома анода,
  и на освободившееся место переходит электрон с L-оболочки (линия $K\alpha$) или M-оболочки ($K\beta$). У меди
  $\lambda_{K\alpha} = 0.15406$~нм (на деле дублет $K\alpha_1/K\alpha_2$), $\lambda_{K\beta} = 0.13922$~нм.
\end{itemize}
Для дифракции нужна одна длина волны, поэтому $K\beta$ гасят \textbf{фильтром}: тонкой фольгой из элемента, у
которого край поглощения лежит между $K\alpha$ и $K\beta$. Для меди это никель: излучение короче 0.1488~нм выбивает
K-электроны никеля и сильно поглощается, а $K\alpha$ проходит почти без потерь.

\fig[0.62]{m1_spectrum}{Спектр трубки с медным анодом при 30~кВ: тормозной фон, начинающийся с $\lambda_{\min}$,
и линии $K\alpha$, $K\beta$ (расчёт по закону Крамерса и гауссовым линиям). Синим --- после никелевого фильтра: $K\beta$
почти погашена, $K\alpha$ ослаблена мало.}

\textbf{Поглощение} следует закону $I = I_0e^{-\mu x}$. Для циркония и излучения Cu $K\alpha$ массовый коэффициент
поглощения около 140~см$^2$/г, то есть $\mu\approx900$~см$^{-1}$: интенсивность падает в $e$ раз на 11~мкм пути. При
съёмке «на отражение» луч входит и выходит под углом, поэтому сигнал собирается с глубины всего нескольких
микрометров. Обычный дифрактометр видит поверхность образца, а не толщу стенки трубы.

\section{Условие Брэгга}

Пусть волна длины $\lambda$ падает под углом $\theta$ на семейство параллельных атомных плоскостей с расстоянием $d$
(рис.~\ref{fig:m2_bragg}а). Волны, отражённые соседними плоскостями, проходят пути, различающиеся на $2d\sin\theta$.
Они усиливают друг друга, только если эта разность --- целое число длин волн:
\begin{keyf}[Условие Брэгга]
\[
n\lambda = 2d\sin\theta .
\]
\end{keyf}
Отражение $n$-го порядка от плоскостей $(hkl)$ удобно считать отражением первого порядка от воображаемых плоскостей
$(nh\,nk\,nl)$ с расстоянием $d/n$. Поэтому обычно пишут просто $\lambda = 2d_{hkl}\sin\theta$. Угол между падающим
и отражённым лучами --- $2\theta$; дифрактограмма --- это интенсивность в зависимости от $2\theta$.

\fig{m2_bragg}{а) Условие Брэгга: волны, отражённые соседними плоскостями, проходят пути, различающиеся на
$2d\sin\theta$ (жёлтые отрезки). б) Обратная решётка и сфера Эвальда: дифракция возникает, когда узел обратной
решётки $\mathbf g_{hkl}$ лежит на сфере радиуса $1/\lambda$, проходящей через начало координат.}

\begin{idea}
Почему отражения такие острые? В кристалле не две, а тысячи плоскостей. При малом отклонении угла от брэгговского
волна от каждой следующей плоскости сдвинута по фазе ещё чуть-чуть, и к тысячной плоскости набегает полный
оборот --- волны гасят друг друга попарно. Чем больше плоскостей участвует, тем уже пик. Поэтому ширина пика
говорит о размере кристаллитов (раздел~\ref{sec:broad}).
\end{idea}

\textbf{Обратная решётка} --- удобная переформулировка. Каждому семейству плоскостей $(hkl)$ сопоставим вектор
$\mathbf g_{hkl}$, перпендикулярный плоскостям, длиной $1/d_{hkl}$. Концы всех таких векторов образуют решётку ---
обратную. Условие Брэгга эквивалентно векторному равенству $\mathbf k - \mathbf k_0 = \mathbf g_{hkl}$, где
$\mathbf k_0$ и $\mathbf k$ --- волновые векторы падающей и рассеянной волн длиной $1/\lambda$. Геометрически: построим
сферу радиуса $1/\lambda$ с центром в начале вектора $\mathbf k_0$ так, чтобы его конец был в начале координат обратной
решётки; отражение есть для каждого узла, попавшего на сферу (\textbf{построение Эвальда},
рис.~\ref{fig:m2_bragg}б). Поворачивая кристалл или беря много кристаллитов разной ориентации (порошок), мы
«протаскиваем» узлы через сферу.

\section{Интенсивность отражений}

Положение пика даёт $d_{hkl}$, то есть геометрию решётки. Интенсивность говорит, какие атомы и где сидят в ячейке.
Каждый атом рассеивает с амплитудой $f_j$ --- \textbf{атомным фактором}, который равен числу электронов при
$\theta = 0$ и падает с ростом $\sin\theta/\lambda$ (электронное облако имеет размер). Волны от атомов ячейки
складываются с фазами, зависящими от их координат $(x_j, y_j, z_j)$:
\begin{keyf}[Структурный фактор]
\[
F_{hkl} = \sum_j f_j\,\exp\bigl[2\pi i\,(hx_j + ky_j + lz_j)\bigr],\qquad I_{hkl}\propto|F_{hkl}|^2 .
\]
\end{keyf}
Отсюда \textbf{правила погасания}.
\begin{itemize}
\item ОЦК, атомы в $(0,0,0)$ и $(\tfrac12,\tfrac12,\tfrac12)$: $F = f\bigl(1 + e^{i\pi(h+k+l)}\bigr)$ --- отражения есть
  только при чётной сумме $h+k+l$: $(110)$, $(200)$, $(211)$\ldots
\item ГЦК: отражения есть, когда $h, k, l$ все чётные или все нечётные: $(111)$, $(200)$, $(220)$, $(311)$\ldots
\item ГПУ, атомы в $(0,0,0)$ и $(\tfrac13,\tfrac23,\tfrac12)$: $F = f\bigl(1 + e^{2\pi i(h/3 + 2k/3 + l/2)}\bigr)$. Нуль,
  если $l$ нечётно и $h+2k$ делится на 3. Поэтому нет отражения $(0001)$, а первое базисное --- $(0002)$.
\end{itemize}
К $|F|^2$ добавляются множители: \textbf{повторяемость} (сколько эквивалентных плоскостей дают один и тот же пик:
6 для $\{10\bar10\}$, 2 для $(0002)$, 12 для $\{10\bar11\}$), \textbf{фактор Лоренца--поляризации}
$(1+\cos^22\theta)/(\sin^2\theta\cos\theta)$, тепловой фактор Дебая--Уоллера $e^{-2B\sin^2\theta/\lambda^2}$ (колебания
атомов размывают решётку, глава~\ref{ch:atoms}) и поглощение.

\fig{m3_xrd}{Расчётные дифрактограммы для Cu $K\alpha$ по формулам этого раздела. а) $\alpha$-Zr, порошок (случайная
ориентация зёрен); подписаны индексы отражений. б) $\alpha$-Zr, съёмка с наружной поверхности трубы: нормаль к
отражающим плоскостям --- радиальное направление; интенсивности умножены на плотность полюсов схематичной текстуры
трубы (глава~\ref{ch:texture}); пунктир --- порошок. в) $\alpha$-Zr с 5\,\% (по объёму) $\delta$-гидрида: слабые
пики гидрида (зелёные) рядом с пиками матрицы.}

На рис.~\ref{fig:m3_xrd} показаны дифрактограммы, посчитанные именно так. У порошка $\alpha$-Zr сильнейший пик ---
$(10\bar11)$ при $2\theta = 36.5^\circ$. С поверхности трубы картина совсем другая: пики призм $(10\bar10)$ и
$(11\bar20)$ почти исчезли, а $(10\bar13)$ и $(10\bar12)$ выросли в разы. Причина --- текстура: оси $c$ большинства
зёрен наклонены к радиусу на 30--40\textdegree, и в отражающее положение попадают плоскости, нормали которых
наклонены к оси $c$ на такой же угол. Сравнение интенсивностей с порошковыми --- простейший способ оценить текстуру
(глава~\ref{ch:texture}).

\begin{zrbox}
Гидрид в цирконии видно на дифрактограмме, но с трудом. Его пики $(111)_\delta$ при 32.4\textdegree{} и
$(200)_\delta$ при 37.6\textdegree{} лежат вплотную к пикам матрицы $(10\bar10)$ и $(10\bar11)$, а доля гидрида при
сотнях ppm водорода --- около процента. В лабораторном дифрактометре это на грани чувствительности. Поэтому
растворение и выпадение гидридов по дифракции изучают на синхротроне: яркий пучок высокой энергии проходит всю
стенку насквозь, и за секунды записывается вся картина (раздел~\ref{sec:synchro}). По сдвигу пиков гидрида при этом
измеряют и деформации в самих гидридах.
\end{zrbox}

\section{Ширина и сдвиг пиков}\label{sec:broad}

\textbf{Размер кристаллитов.} Если кристаллит содержит мало плоскостей, пик уширяется. Формула Шеррера:
\begin{equation}
\beta = \frac{K\lambda}{L\cos\theta},\qquad K\approx0.9,
\end{equation}
где $\beta$ --- уширение (в радианах шкалы $2\theta$), $L$ --- размер кристаллита по нормали к плоскостям. При $L = 1$~мкм
уширение ничтожно, при 50~нм --- около 0.2\textdegree, при 15~нм --- около 0.6\textdegree{} (рис.~\ref{fig:m4_peaks}а).
Поэтому так оценивают только наноразмерные объекты; зёрна оболочки в несколько микрометров «на ширине» не видны.

\textbf{Микродеформации.} Неоднородные деформации (дислокации, поля вокруг выделений) тоже уширяют пики, но по-другому:
$\beta\propto\operatorname{tg}\theta$. Разделить два вклада позволяет построение Уильямсона--Холла: $\beta\cos\theta$ против
$\sin\theta$ --- отрезок на оси даёт размер, наклон --- микродеформацию. Из ширины пиков холоднодеформированного металла
так оценивают плотность дислокаций.

\textbf{Макронапряжения: метод $\sin^2\psi$.} Однородное упругое напряжение \emph{сдвигает} пики. Пусть в поверхностном
слое действует напряжение $\sigma$ вдоль некоторого направления. Наклоним образец на угол $\psi$ в плоскости этого
направления и будем мерить $d$ для одного и того же отражения. Плоскости с нормалью, наклонённой на $\psi$ от нормали к
поверхности, растянуты на
\begin{equation}
\eps_\psi = \frac{d_\psi - d_0}{d_0} = \frac{1+\nu}{E}\,\sigma\sin^2\psi - \frac{\nu}{E}(\sigma_{11}+\sigma_{22}).
\end{equation}
Зависимость $d$ от $\sin^2\psi$ --- прямая, наклон которой пропорционален напряжению (рис.~\ref{fig:m4_peaks}б).
Для циркония вместо изотропных $E$ и $\nu$ берут «рентгеновские упругие постоянные» данного отражения: в разных
зёрнах решётка нагружена по-разному.

\fig{m4_peaks}{а) Уширение пика $(10\bar11)$ $\alpha$-Zr по Шерреру для кристаллитов 1~мкм, 50 и 15~нм. б) Метод
$\sin^2\psi$: межплоскостное расстояние $d_{10\bar11}$ от наклона образца при одноосном напряжении 0 и $\pm200$~МПа
($E = 99$~ГПа, $\nu = 0.34$).}

\begin{warnbox}
Дифракция меряет деформацию решётки только в тех зёрнах, которые стоят в отражающем положении. В анизотропном
поликристалле эти зёрна нагружены не так, как металл в среднем: межзёренные напряжения второго рода
(глава~\ref{ch:plasticity}) у разных отражений разные. Поэтому напряжения по разным пикам циркония расходятся, и
это не ошибка, а информация.
\end{warnbox}

\section{Синхротрон и нейтроны}\label{sec:synchro}

\textbf{Синхротронное излучение} --- рентген от электронов, движущихся по кругу в ускорителе. Оно на много порядков ярче
трубки, а энергию кванта можно выбрать высокой (80--100~кэВ). Такое жёсткое излучение проходит миллиметры циркония
насквозь, и дифракцию снимают «на просвет» двумерным детектором: в одном кадре --- кольца всех отражений, а их
неоднородность по азимуту сразу показывает текстуру. Кадр занимает доли секунды, поэтому можно греть и охлаждать
образец под нагрузкой и следить, при какой температуре появляется и исчезает гидрид, какая у него деформация и как
нагрузка перераспределяется между фазами. Многие современные данные о растворимости, гистерезисе и выпадении под
напряжением получены именно так.

\textbf{Нейтроны} рассеиваются на ядрах, а не на электронах, поэтому видят лёгкие атомы не хуже тяжёлых. Водород даёт
огромное некогерентное рассеяние --- по нему меряют содержание водорода и его распределение (нейтронная
радиография труб). Для дифракции берут дейтерий, у которого некогерентное рассеяние мало. Нейтроны проникают на
сантиметры, поэтому нейтронная дифракция меряет текстуру и напряжения в объёме, а не на поверхности.

\section{Электронная микроскопия}

Электроны с энергией в десятки--сотни кэВ имеют длину волны в тысячные доли нанометра. Их можно фокусировать
магнитными линзами и получать изображения с разрешением вплоть до атомного.

\textbf{Растровый электронный микроскоп (РЭМ).} Тонкий пучок обегает поверхность построчно. Вторичные электроны (с
поверхности) дают рельеф, отражённые (обратнорассеянные) --- контраст по среднему атомному номеру: гидрид и
матрица различаются слабо, а интерметаллиды с железом и ниобием видны хорошо. Характеристическое рентгеновское
излучение (энергодисперсионный анализ, ЭДС) даёт химический состав в точке размером около микрометра --- но
водород этим методом не виден.

\textbf{Дифракция обратнорассеянных электронов (EBSD).} Образец наклоняют на 70\textdegree; электроны, рассеянные в
приповерхностном слое, дифрагируют на плоскостях решётки и дают на экране \textbf{линии Кикучи} --- полосы, каждая из
которых отвечает семейству плоскостей. По их рисунку программа определяет ориентацию кристалла в данной точке.
Обегая поверхность с шагом 0.05--0.5~мкм, получают \textbf{карту ориентаций}: зёрна, их границы и разориентации,
текстуру по тысячам зёрен, ориентационное соотношение гидрида с матрицей и то, к какой плоскости матрицы он
прилегает. Для EBSD нужна очень чистая, без наклёпа, поверхность (электрополировка или ионная полировка).

\textbf{Просвечивающий электронный микроскоп (ПЭМ).} Пучок проходит через фольгу толщиной около 100~нм. Видны отдельные
дислокации, дефекты упаковки, микропластинки гидрида и их когерентность; электронограммы дают ориентационные
соотношения. Цена --- крошечный объём исследования и сложная подготовка (а при утонении фольги гидриды могут
выпадать или растворяться). Последовательная резка ионным пучком в РЭМ (FIB-SEM) даёт трёхмерную реконструкцию
пачек гидридов.

\section{Оптическая металлография}

Самый распространённый способ увидеть гидриды --- шлиф. Образец шлифуют, полируют и \textbf{травят}: травитель
(обычно смесь плавиковой и азотной кислот с водой) растворяет гидрид и матрицу с разной скоростью, и на месте
гидридов остаются канавки, которые в светлом поле выглядят тёмными линиями.

Предел разрешения светового микроскопа $\delta\approx0.61\lambda/\mathrm{NA}$: при $\lambda = 0.55$~мкм и числовой апертуре
0.5 это около 0.7~мкм. Микропластинка толщиной 0.1--0.3~мкм оптически не разрешается --- без травления она дала бы
бледную размытую полоску шириной около микрометра. Травление делает её видимой, но расширяет след на 1--2~мкм с
каждой стороны (рис.~\ref{fig:m6_optics}а).

\fig{m6_optics}{а) Профиль затемнения поперёк пластинки гидрида толщиной 0.3~мкм: сама пластинка, её изображение
в оптике (дифракционное размытие) и изображения канавок после травления, расширяющего след на 1 и 2~мкм с каждой
стороны. б) Видимая доля площади, занятая линиями гидрида, от утолщения травлением при густоте линий 10 и
15~мм/мм$^2$ (длина линий на единицу площади).}

\begin{zrbox}
Следствия для сравнения модели со шлифами.
\begin{itemize}
\item \textbf{Толщину} гидридов по травленому шлифу мерить нельзя --- видна ширина канавки, а не пластинки. Видимая
  доля площади в несколько раз больше истинной доли гидрида (рис.~\ref{fig:m6_optics}б).
\item \textbf{Длина и ориентация} линий травлением искажаются мало --- поэтому именно они (доля радиальных гидридов,
  распределение длин, густота линий) служат основными мерами.
\item Близкие пластинки на расстоянии меньше 2--4~мкм сливаются в одну линию. Пачка микропластинок на шлифе ---
  одна «макропластинка», и сравнивать с ней нужно не отдельные пластинки модели, а их проекцию с тем же
  утолщением.
\item Результат зависит от оператора: времени травления, свежести травителя, порога при обработке изображения.
  Сопоставимы только шлифы, приготовленные одинаково.
\end{itemize}
\end{zrbox}

Ориентацию гидридов на шлифе описывают числами: долей радиальных гидридов (RHF --- суммарная длина линий,
отклонённых от окружного направления больше чем на заданный угол, обычно 40--45\textdegree, к общей длине) и родственными
мерами вроде $F_n$ (доля пластинок в заданном угловом секторе). Их считают по изображению: порог, выделение линий,
скелетизация, разбиение на отрезки и подсчёт углов.

\section{Тепловые методы}

\textbf{Дифференциальная сканирующая калориметрия (ДСК).} Образец и эталон нагревают или охлаждают с постоянной
скоростью и измеряют разность тепловых потоков. Плавный фон --- теплоёмкость, пики --- превращения. Растворение
гидрида при нагреве поглощает теплоту (эндотермический пик), выпадение при охлаждении её выделяет
(экзотермический).

\fig[0.68]{m5_dsc}{Расчётные кривые ДСК для Э635 со 160~ppm водорода (10~К/мин, инерция прибора 0.5~мин): тепловой
поток пропорционален скорости растворения или выпадения, $q\propto|\dd C/\dd T|$, по линиям TSSD и TSSP главы~\ref{ch:solutions}.
При нагреве пик растворения заканчивается на TSSD; при охлаждении пик выпадения начинается на TSSP.}

Как читать кривые (рис.~\ref{fig:m5_dsc}). При нагреве гидрид растворяется всё быстрее (растворимость растёт
экспоненциально), и пик обрывается, когда растворился последний гидрид, --- при температуре TSSD для данного
содержания водорода. За TSSD берут температуру конца пика (часто --- его максимума или точки отрыва от фона; от этого
выбора и от скорости нагрева значения у разных авторов расходятся на 5--15~К). При охлаждении ничего не происходит
до TSSP, затем выпадение начинается резко --- по началу экзотермического пика. Серия образцов с разным водородом
даёт линии растворимости.

\textbf{Дилатометрия} меряет длину образца при нагреве: гидрид объёмнее матрицы, и его растворение даёт излом на
кривой теплового расширения. Тот же метод фиксирует превращение $\alpha\to\beta$. Применяют также
\textbf{электросопротивление} (растворённый водород рассеивает электроны сильнее, чем гидрид той же массы) и
\textbf{внутреннее трение}.

\section{Измерение содержания водорода}

Стандартный метод --- \textbf{экстракция плавлением в инертном газе}: пробу массой 0.1--0.5~г расплавляют в графитовом
тигле в токе аргона, выделившийся водород уносится газом и измеряется детектором теплопроводности. Точность ---
несколько ppm или несколько процентов от измеряемой величины. Метод разрушающий и даёт среднее по пробе, поэтому
распределение водорода по стенке или длине трубы меряют, разрезая образец на части. Неразрушающие способы ---
нейтронная радиография и ДСК (по площади пика или по температуре TSSD).

\begin{zrbox}
Образцы для опытов по переориентации наводороживают из газа или электролитически, затем отжигают для выравнивания
водорода. Неравномерность водорода по стенке (поверхностный гидридный слой, больше у наружной или внутренней
поверхности) прямо влияет на картину гидридов. Поэтому в моделях и сравнениях важно знать не только среднее
содержание, но и как его меряли.
\end{zrbox}

\begin{exercise}
Найдите угол $2\theta$ отражения $(0002)$ $\alpha$-Zr для Cu $K\alpha$ ($c = 0.5147$~нм). А для молибденового излучения
($\lambda = 0.0709$~нм)?
\end{exercise}

\begin{exercise}
Вычислите структурный фактор ГПУ для отражений $(0001)$, $(0002)$, $(10\bar10)$ и $(10\bar11)$. Какие из них погашены?
\end{exercise}

\begin{exercise}
Оцените по Шерреру уширение пика $(10\bar11)$ при $L = 30$~нм. Сможете ли вы различить его на приборе с
инструментальной шириной 0.08\textdegree?
\end{exercise}

\begin{exercise}
Почему температуру TSSD определяют по нагреву, а TSSP --- по охлаждению? Что будет с измеренной TSSD, если
скорость нагрева увеличить?
\end{exercise}
